\chapter{Leitfaden für Mitglieder}
\label{leitfaden-mitglieder}
Dieses Kapitel beschreibt in \autoref{mitglieder-anleitungen}, wie alltägliche Aufgaben erledigt und Dokumente erstellt und kompiliert (also die PDF-Dateien erstellt) werden, und in \autoref{mitglieder-referenz}, welche Vorlagen es gibt und was sie festlegen.

\section{Anleitungen}
\label{mitglieder-anleitungen}
\subsection{Verfassen eines Sitzungsprotokolls}
\label{protokoll}
Um die Beschlüsse aus den Sitzungen des Fachschaftsrats festzuhalten, werden \textit{Sitzungsprotokolle} angefertigt. Dafür gibt es eine eigens erstellte \hyperref[latex-fsrprotocol]{\LaTeX-Dokumentenklasse}, die so gestaltet ist, dass auch Mitglieder ohne \LaTeX-Erfahrung damit arbeiten können.

Grundlagen zu \LaTeX{}, zum Kompilieren der PDF-Datei und zum Umgang mit Fehlermeldungen finden sich in \autoref{latex-uebersicht}, \autoref{latex-kompilieren} und \autoref{latex-fehler-finden}.

\subsubsection{Erstellen eines neuen Protokolls}
\label{protokoll-erstellen}
Auf dem Arbeitsrechner des Fachschaftsrats ist eine Vorlage für Sitzungsprotokolle hinterlegt. Mit ihr lässt sich ein neues Protokoll direkt im Dateimanager anlegen:

\notebox{Die folgende Anleitung gilt für den eingerichteten Arbeitsrechner des Fachschaftsrats. Auf eigenen Rechnern müssen zunächst die in \autoref{latex-eigener-rechner} beschriebenen Voraussetzungen erfüllt sein. Statt über das Kontextmenü wird dort eine leere \texttt{.tex}-Datei angelegt und mit den Codeblöcken aus \autoref{protokoll-schreiben} gefüllt.}

\begin{enumerate}
	% TODO: Nach Fertigstellung des Servers den Pfad zu Dokumente/Protokolle/JJJJ ändern.
	\item Im Dateimanager den Ordner \texttt{Dokumente/webspace/Protokolle/JJJJ} für das Jahr der Sitzung öffnen. Falls der Ordner noch nicht existiert, muss er zuerst angelegt werden.
	\item Mit einem Rechtsklick auf eine freie Fläche das Kontextmenü öffnen und
	      \textit{Neu erstellen} \textrightarrow{} \texttt{fsrprotocol.tex} auswählen.
	\item Die neue Datei nach dem Schema \texttt{JJJJ-MM-DD\_Sitzungsprotokoll\_Name.tex} umbenennen,
	      z.\,B. \texttt{2026-10-14\_Sitzungsprotokoll\_Max-Mustermann.tex}.
	\item Die Datei in einem Editor öffnen (auf dem Arbeitsrechner des Fachschaftsrats: VSCodium) und wie in \autoref{protokoll-schreiben} beschrieben bearbeiten.
\end{enumerate}

\subsubsection{Schreiben eines Protokolls}
\label{protokoll-schreiben}
Im ersten Teil des \LaTeX-Dokuments werden die Dokumentenklasse geladen und die Metadaten der Sitzung festgelegt.

\importantbox{Die Angaben zur Sitzung sind Pflicht: Fehlt eine davon, verweigert \LaTeX{} die \hyperref[latex-kompilieren]{Kompilierung}. Die dabei angezeigte Fehlermeldung ist in \autoref{latex-fehlermeldungen} beschrieben.}

\notebox{Text hinter einem \texttt{\%} ist ein Kommentar und erscheint nicht im PDF. Ein \texttt{\%} im Text wird als \mintinline{latex}{\%} geschrieben (siehe \autoref{latex-sonderzeichen}).}
\begin{Codeblock}[title={Anfang des Dokuments}, unbreakable]{latex}
	\documentclass{fsrprotocol}

	% --- Setze Informationen zu dem Dokument --- %
	% Hier stehen Informationen zur Sitzung. Der Protokollführer muss diese ausfüllen.
	\newcommand{\Sitzungsdatum}{dd.mm.yyyy}
	\newcommand{\Sitzungszeitraum}{00:00 - 23:59 Uhr}
	\newcommand{\Sitzungsort}{Hier Veranstaltungsort eintragen}
	\newcommand{\Sitzungsleitung}{Vorname Nachname}
	\newcommand{\Protokollfuehrer}{Vorname Nachname}
	\newcommand{\Anwesenheitsliste}{Vorname1 Nachname1, Vorname2 Nachname2, ...}
	\newcommand{\Gaesteliste}{Vorname1 Nachname1, Vorname2 Nachname2, ...}

\begin{document}
\end{Codeblock}

Nur die Platzhalter in den geschweiften Klammern der \mintinline{latex}{\newcommand}-Zeilen werden ersetzt; alles andere bleibt unverändert. Waren keine Gäste anwesend, wird bei \mintinline{latex}{\Gaesteliste} \enquote{keine} eingetragen. Hintergründe zu \mintinline{latex}{\documentclass} stehen in \autoref{latex-uebersicht}.

Als Nächstes werden die Tagesordnungspunkte in der Umgebung \texttt{Toplist} aufgelistet:

\begin{Codeblock}[title={Liste der Tagesordnungspunkte}]{latex}
	\begin{Toplist}
		\item Erster Tagesordnungspunkt
		\item Zweiter Tagesordnungspunkt
		\item ...
	\end{Toplist}
\end{Codeblock}

Danach wird jeder Tagesordnungspunkt in einer eigenen \texttt{Top}-Umgebung beschrieben:

\begin{Codeblock}[title={Erstellen eines Tagesordnungspunkts mit Abstimmungen}]{latex}
	\begin{Top}{Titel des ersten Tagesordnungspunkts}
		Hier die ausführliche Beschreibung und Diskussion zum ersten TOP.
		% Abstimmungen sind optional
		\Abstimmung[]{Erste Abstimmung zum ersten Tagesordnungspunkt}{2}{1}{2}
		\Abstimmung[vertagt]{Zweite Abstimmung zum ersten Tagesordnungspunkt}{}{}{}
	\end{Top}
\end{Codeblock}

Innerhalb eines Tagesordnungspunkts können eine oder mehrere Abstimmungen festgehalten werden. Dazu dient der Befehl \mintinline{latex}{\Abstimmung[Vermerk]{Gegenstand}{Ja}{Nein}{Enthaltungen}}. Bei einer regulären Abstimmung bleiben die eckigen Klammern leer und alle drei Stimmfelder werden mit einer Zahl gefüllt, auch wenn diese 0 ist.

Wurde eine Abstimmung nicht durchgeführt, etwa weil sie vertagt wurde, wird statt der Stimmen ein Vermerk eingetragen: Der Vermerk steht in den eckigen Klammern, die drei Stimmfelder bleiben leer (siehe die zweite Abstimmung im Beispiel). Wie das Ergebnis einer Abstimmung bestimmt wird, ist in \autoref{latex-fsrprotocol} beschrieben.

Der Gegenstand einer Abstimmung sollte möglichst im Wortlaut des Antrags festgehalten werden, da das Protokoll als Nachweis für die Beschlüsse dient.

Zum Schluss muss das Dokument durch Beenden der \texttt{document}-Umgebung abgeschlossen werden:

\begin{Codeblock}[title={Abschließen des Dokuments}]{latex}
	\end{document}
\end{Codeblock}

\subsubsection{Abschließen und Ablegen}
\label{protokoll-ablegen}
\begin{enumerate}
	\item Das fertige Protokoll \hyperref[latex-kompilieren]{kompilieren}.
	\item Die erzeugte PDF-Datei auf Vollständigkeit und Tippfehler prüfen und gegebenenfalls korrigieren und neu kompilieren.
	\item Die PDF-Datei aus dem Ordner \texttt{build/} in den Jahresordner verschieben, sodass \texttt{.tex}- und \texttt{.pdf}-Datei zusammen abgelegt sind.
\end{enumerate}

Bei Unklarheiten zur Erstellung eines Protokolls hilft der aktuelle Administrator weiter. Auch die \texttt{.tex}-Dateien älterer Protokolle können als Orientierung hilfreich sein.

\subsection{Kompilieren einer PDF-Datei}
\label{latex-kompilieren}
Auf dem Arbeitsrechner des Fachschaftsrats werden Dokumente in VSCodium mit der Erweiterung LaTeX Workshop kompiliert, die dort bereits für LuaLaTeX eingerichtet ist:

\begin{enumerate}
	\item Die \texttt{.tex}-Datei in VSCodium öffnen.
	\item Das Dokument mit \texttt{Strg+Alt+B} oder über den grünen Pfeil \textit{Build LaTeX project} oben rechts kompilieren. Standardmäßig wird das Dokument außerdem bei jedem Speichern automatisch neu kompiliert.
	\item Die erzeugte PDF-Datei mit \texttt{Strg+Alt+V} oder über das Symbol \textit{View LaTeX PDF} oben rechts in einer Vorschau neben dem Quelltext öffnen. Die Vorschau aktualisiert sich nach jedem Kompilieren automatisch.
\end{enumerate}
% TODO: Prüfen, ob Tastenkürzel und automatisches Kompilieren beim Speichern auf tesla so eingerichtet sind.

Die PDF-Datei wird im automatisch erstellten Ordner \texttt{build/} abgelegt. Die fertige PDF-Datei kann anschließend beliebig verschoben werden.

\subsection{Erstellen von Dokumenten auf dem eigenen Rechner}
\label{latex-eigener-rechner}
Dokumente können auch auf einem eigenen Rechner kompiliert werden. Dort muss ebenfalls LuaLaTeX verwendet werden (z.\,B. mit dem Terminal-Befehl \texttt{latexmk -lualatex datei.tex}). Damit \LaTeX{} die Vorlagen des Fachschaftsrats findet, werden sie wie folgt eingerichtet:
\begin{enumerate}
	\item Die Vorlagen aus dem \href{https://github.com/fsrphysik-uds/latex-templates}{Repository} herunterladen (\textit{Code} \textrightarrow{} \textit{Download ZIP}) und entpacken.
	\item Folgende Dateien in denselben Ordner wie die \texttt{.tex}-Datei kopieren:
	      \begin{itemize}
		      \item das Basispaket \texttt{fsrbase.sty},
		      \item die jeweils verwendete Dokumentenklasse, z.\,B. \texttt{fsrprotocol.cls},
		      \item das Logo \texttt{media/fsr-logo.png}, falls der Befehl \mintinline{latex}{\fsrlogo} verwendet wird (bei Sitzungsprotokollen immer, siehe \autoref{latex-fsrprotocol}).
	      \end{itemize}
\end{enumerate}

\notebox{Alternativ können die Dateien auch systemweit installiert werden, sodass sie in jedem Ordner zur Verfügung stehen und nicht in jeden Arbeitsordner kopiert werden müssen. Wie das funktioniert, hängt von der installierten \LaTeX{}-Distribution ab (z.\,B. TeX~Live oder MiKTeX). Details beschreibt die Dokumentation der jeweiligen Distribution.}

Bei Online-Editoren wie Overleaf ist eine systemweite Installation nicht möglich. Dort müssen die genannten Dateien immer in das jeweilige Projekt hochgeladen werden und neben der \texttt{.tex}-Datei liegen. Außerdem muss in den Projekteinstellungen (\textit{Menü} \textrightarrow{} \textit{Compiler}) LuaLaTeX als Compiler ausgewählt werden.

\subsection{Finden von Fehlern}
\label{latex-fehler-finden}
Fehler werden in LaTeX Workshop im Reiter \textit{Probleme} am unteren Bildschirmrand angezeigt.

Bei Fehlern, die nicht in \autoref{latex-fehlermeldungen} beschrieben sind, hilft folgendes Vorgehen:
\begin{enumerate}
	\item Im Reiter \textit{Probleme} auf die Fehlermeldung klicken. Der Editor springt dann zu der Zeile, in der \LaTeX{} den Fehler gefunden hat. Die eigentliche Ursache liegt oft in dieser oder in der Zeile davor, z.\,B. eine vergessene Klammer.
	\item Lässt sich der Fehler so nicht finden, die letzten Änderungen rückgängig machen (\texttt{Strg+Z}), bis das Dokument wieder kompiliert, und die Änderung erneut vornehmen.
	\item Hilft das nicht weiter, kann der aktuelle Administrator weiterhelfen.
\end{enumerate}

\subsection{Erstellen von Dokumenten ohne passende Dokumentenklasse}
\label{latex-ohne-dokumentenklasse}
Für Dokumente, für die es noch keine eigene Dokumentenklasse gibt (z.\,B. Aushänge oder Schreiben), wird das \hyperref[latex-fsrbase]{Basispaket} direkt eingebunden:
\begin{enumerate}
	\item Die passende \KOMAScript{}-Klasse aus \autoref{latex-koma} mit \mintinline{latex}{\documentclass} laden, z.\,B. \texttt{scrartcl} für einen Aushang.
	\item Direkt darunter das Basispaket mit \mintinline{latex}{\usepackage{fsrbase}} einbinden.
	\item Den Inhalt zwischen \mintinline{latex}{\begin{document}} und \mintinline{latex}{\end{document}} schreiben.
	\item Das Dokument \hyperref[latex-kompilieren]{kompilieren}.
\end{enumerate}

\begin{Codeblock}[title={Minimales Dokument mit dem Basispaket}]{latex}
\documentclass{scrartcl}
\usepackage{fsrbase}

\begin{document}
\fsrlogo[4cm]

\section*{Überschrift}
Hier steht der Inhalt des Dokuments.
\end{document}
\end{Codeblock}


\section{Referenz}
\label{mitglieder-referenz}
\subsection{Übersicht über Paket und Dokumentenklassen}
\label{latex-uebersicht}
Die \LaTeX{}-Vorlagen des Fachschaftsrats bestehen aus dem Basispaket \texttt{fsrbase} und darauf aufbauenden Dokumentenklassen. Eine Dokumentenklasse legt Layout, Gestaltung und wiederkehrende Elemente eines Dokumenttyps fest, z.\,B. bei \hyperref[protokoll]{Sitzungsprotokollen}.

\begin{center}
	\begin{tabularx}{\linewidth}{lX}
		\toprule
		\textbf{Name}               & \textbf{Zweck}                                                                                                                   \\
		\midrule
		\texttt{fsrbase}            & \hyperref[latex-fsrbase]{Basispaket}: legt Schriftarten, Farben, deutsche Sprache und Anführungszeichen fest und stellt den Befehl \texttt{\textbackslash fsrlogo} bereit. \\
		\texttt{fsrprotocol}        & \hyperref[latex-fsrprotocol]{Dokumentenklasse} für \hyperref[protokoll]{Sitzungsprotokolle}.                                                   \\
		\bottomrule
	\end{tabularx}
\end{center}
% TODO: Weitere Dokumentenklassen ergänzen, sobald diese existieren.

Welche Dokumentenklasse verwendet wird, legt der Befehl \mintinline{latex}{\documentclass} in der ersten Zeile des Dokuments fest. Ein Sitzungsprotokoll beginnt beispielsweise mit \mintinline{latex}{\documentclass{fsrprotocol}}. Das Basispaket \texttt{fsrbase} wird dabei bei allen Dokumentenklassen des Fachschaftsrats automatisch mitgeladen.

\subsubsection{\KOMAScript{}-Klassen}
\label{latex-koma}
Dokumente des Fachschaftsrats verwenden statt der Standardklassen von \LaTeX{} die entsprechenden \KOMAScript{}-Klassen. Auch \texttt{fsrprotocol} baut auf \texttt{scrartcl} auf. Für Dokumenttypen ohne eigene Dokumentenklasse wird das Basispaket direkt zusammen mit einer dieser Klassen geladen (siehe \autoref{latex-ohne-dokumentenklasse}).

\begin{center}
	\begin{tabularx}{\linewidth}{llX}
		\toprule
		\textbf{Standardklasse} & \textbf{\KOMAScript{}-Klasse} & \textbf{Typische Verwendung}          \\
		\midrule
		\texttt{article}        & \texttt{scrartcl}             & Kurze Dokumente, z.\,B. Aushänge      \\
		\texttt{report}         & \texttt{scrreprt}             & Längere Dokumente mit Kapiteln        \\
		\texttt{book}           & \texttt{scrbook}              & Bücher, doppelseitiger Satz           \\
		\texttt{letter}         & \texttt{scrlttr2}             & Briefe und Schreiben                  \\
		\bottomrule
	\end{tabularx}
\end{center}

\subsection{Das Basispaket \texttt{fsrbase}}
\label{latex-fsrbase}
Das Basispaket sorgt dafür, dass alle Dokumente des Fachschaftsrats einheitlich aussehen. Es legt Folgendes fest:
\begin{itemize}
	\item \textbf{Kompilierprogramm:} Das Paket setzt LuaLaTeX voraus und bricht mit einer \hyperref[latex-fehlermeldungen]{Fehlermeldung} ab, wenn ein anderes Programm (z.\,B. pdfLaTeX) verwendet wird.
	\item \textbf{Schriftarten:} TeX Gyre Termes als Serifenschrift und TeX Gyre Heros als serifenlose Schrift (über \texttt{fontspec}).
	\item \textbf{Sprache:} Deutsche Silbentrennung und Typografie über \texttt{babel} mit der Option \texttt{ngerman}.
	\item \textbf{Anführungszeichen:} Über \texttt{csquotes} werden gerade Anführungszeichen \texttt{\string"...\string"} im Text automatisch zu deutschen Anführungszeichen "..." umgewandelt. Alternativ kann \mintinline{latex}{\enquote{...}} verwendet werden.
	\item \textbf{Weitere Pakete:} \texttt{microtype} (bessere Typografie/Blocksatz), \texttt{xcolor} (Farben) und \texttt{graphicx} (Einbinden von Bildern) werden geladen und müssen nicht erneut eingebunden werden.
\end{itemize}

Für den Inhalt stellt das Basispaket folgende Farben und Befehle bereit:
\begin{description}
	\item[\texttt{fsrcolor}] Akzentfarbe des Fachschaftsrats (Schwarz, RGB 0,\,0,\,0).
	\item[\texttt{fsrgrey}] Grauton des Fachschaftsrats (HTML \texttt{4D4D4D}).
	\item[{\texttt{\textbackslash fsrlogo[Höhe]}}] Setzt das Fachschaftslogo aus der Datei \texttt{fsr-logo.png} in der angegebenen Höhe. Ohne Angabe der Höhe wird das Logo 5\,cm hoch gesetzt.
\end{description}

Die Farben werden wie andere \texttt{xcolor}-Farben verwendet, z.\,B. mit \mintinline{latex}{\textcolor{fsrgrey}{Text}}.

\subsection{Die Dokumentenklasse \texttt{fsrprotocol}}
\label{latex-fsrprotocol}
Die Dokumentenklasse für \hyperref[protokoll]{Sitzungsprotokolle} übernimmt Layout und Gestaltung sowie die Nummerierung der Tagesordnungspunkte und Abstimmungen. Sie legt Folgendes fest:
\begin{itemize}
	\item \textbf{Grundlage:} \KOMAScript{}-Klasse \texttt{scrartcl} mit Schriftgröße 11\,pt, Papierformat A4, halbzeiligem Absatzabstand statt Einrückung (\texttt{parskip=half}) und Satzspiegel \texttt{DIV=12}. Das \hyperref[latex-fsrbase]{Basispaket} wird automatisch geladen.
	\item \textbf{Kopf:} Titel \enquote{Sitzungsprotokoll}, Logo in 4,3\,cm Höhe (über \mintinline{latex}{\fsrlogo}, daher wird immer die Datei \texttt{fsr-logo.png} benötigt), Angaben zum Veranstaltungsort und der Sitzung aus der folgenden Tabelle.
\end{itemize}

Die folgenden Befehle legen die Informationen zur jeweiligen Sitzung fest. Sie werden mit \texttt{\textbackslash newcommand} vor \texttt{\textbackslash begin\{document\}} festgelegt. Die Angaben zur Sitzung sind Pflicht: Fehlt eine davon, verweigert \LaTeX{} die Kompilierung (siehe \autoref{latex-fehlermeldungen}). Waren keine Gäste anwesend, wird \texttt{\textbackslash Gaesteliste} auf \enquote{keine} gesetzt.

\begin{center}
	\begin{tabularx}{\linewidth}{llX}
		\toprule
		\textbf{Befehl}                           & \textbf{Im Kopf als} & \textbf{Platzhalter in der Vorlage}         \\
		\midrule
		\texttt{\textbackslash Sitzungsdatum}     & Datum                & dd.mm.yyyy                                  \\
		\texttt{\textbackslash Sitzungszeitraum}  & Zeitraum             & 00:00 - 23:59 Uhr                           \\
		\texttt{\textbackslash Sitzungsort}       & Ort                  & Hier Veranstaltungsort eintragen            \\
		\texttt{\textbackslash Sitzungsleitung}   & Sitzungsleitung      & Vorname Nachname                            \\
		\texttt{\textbackslash Protokollfuehrer}  & Protokollführer      & Vorname Nachname                            \\
		\texttt{\textbackslash Anwesenheitsliste} & Anwesende            & Vorname1 Nachname1, Vorname2 Nachname2, ... \\
		\texttt{\textbackslash Gaesteliste}       & Gäste                & Vorname1 Nachname1, Vorname2 Nachname2, ... \\
		\bottomrule
	\end{tabularx}
\end{center}

Für den Inhalt des Protokolls stellt die Dokumentenklasse folgende Umgebungen und Befehle bereit:
\begin{description}
	\item[\texttt{Toplist}] Umgebung, die die Tagesordnungspunkte automatisch als Aufzählung darstellt.
	\item[\texttt{Top}] Umgebung für die ausführliche Beschreibung eines Tagesordnungspunkts. Ein Tagesordnungspunkt wird mit \texttt{\textbackslash begin\{Top\}\{Titel\}} begonnen und mit \texttt{\textbackslash end\{Top\}} beendet.
	\item[{\texttt{\textbackslash Abstimmung[Vermerk]\{Gegenstand\}\{Ja\}\{Nein\}\{Enthaltungen\}}}] Hält innerhalb eines Tagesordnungspunkts eine Abstimmung fest. Bei einer regulären Abstimmung bleibt der Vermerk leer und alle drei Stimmfelder enthalten eine Zahl. Fehlt bei den Ja- oder Nein-Stimmen eine Zahl, bricht die Kompilierung mit einem Fehler ab. Das Ergebnis wird automatisch bestimmt: Gibt es mehr Ja- als Nein-Stimmen, wird die Abstimmung als "angenommen" vermerkt, ansonsten als "abgelehnt". Wurde eine Abstimmung nicht durchgeführt, wird statt der Stimmen der Vermerk eingetragen, die drei Stimmfelder bleiben dann leer.
\end{description}

\subsection{Sonderzeichen}
\label{latex-sonderzeichen}
Einige Zeichen haben in \LaTeX{} eine besondere Bedeutung, z.\,B. leitet \texttt{\%} einen Kommentar ein. Sie können deshalb nicht einfach so im Text verwendet werden. Sollen sie im Text erscheinen, wird die Eingabe aus der Tabelle verwendet. Bei den meisten Zeichen ist das ein vorangestellter Backslash, z.\,B. \mintinline{latex}{50\%} oder \mintinline{latex}{Q\&A}.

\begin{center}
	\begin{tabularx}{\linewidth}{clX}
		\toprule
		\textbf{Zeichen} & \textbf{Eingabe}            & \textbf{Bedeutung ohne Backslash}        \\
		\midrule
		\texttt{\%}      & \texttt{\textbackslash\%}   & Leitet einen Kommentar ein               \\
		\texttt{\&}      & \texttt{\textbackslash\&}   & Trennt Spalten in Tabellen               \\
		\texttt{\_}      & \texttt{\textbackslash\_}   & Tiefstellung in Formeln                  \\
		\texttt{\#}      & \texttt{\textbackslash\#}   & Platzhalter für Argumente von Befehlen   \\
		\texttt{\$}      & \texttt{\textbackslash\$}   & Beginnt und beendet eine Formel          \\
		\texttt{\{}      & \texttt{\textbackslash\{}   & Beginnt eine Gruppe oder ein Argument    \\
		\texttt{\}}      & \texttt{\textbackslash\}}   & Beendet eine Gruppe oder ein Argument    \\
		\texttt{\textasciitilde} & \texttt{\textbackslash textasciitilde} & Geschütztes Leerzeichen (kein Zeilenumbruch) \\
		\texttt{\textasciicircum} & \texttt{\textbackslash textasciicircum} & Hochstellung in Formeln \\
		\texttt{\textbackslash} & \texttt{\textbackslash textbackslash} & Leitet einen Befehl ein \\
		\bottomrule
	\end{tabularx}
\end{center}

\subsection{Fehlermeldungen der Vorlagen}
\label{latex-fehlermeldungen}
Neben den üblichen \LaTeX{}-Fehlern können die Vorlagen des Fachschaftsrats eigene Fehlermeldungen ausgeben:
\begin{description}
	\item[\texttt{LuaTeX is required to compile this document}] Das Dokument wurde nicht mit LuaLaTeX kompiliert (siehe \autoref{latex-kompilieren} bzw. \autoref{latex-eigener-rechner}).
	\item[\texttt{Makro ... ist nicht definiert}] In einem Sitzungsprotokoll fehlt eine der Pflichtangaben am Anfang des Dokuments (siehe \autoref{protokoll-schreiben}).
	\item[\texttt{Missing number, treated as zero}] In einem Sitzungsprotokoll fehlt bei einer \texttt{\textbackslash Abstimmung} ohne Vermerk die Zahl der Ja- oder Nein-Stimmen (siehe \autoref{latex-fsrprotocol}). Diese Meldung ist ein allgemeiner \LaTeX{}-Fehler und kann auch andere Ursachen haben.
\end{description}

Zum Vorgehen bei anderen Fehlern siehe \autoref{latex-fehler-finden}.

